% !TEX root = ../../tree_notes.tex
% =====================================================================
%  7.2  Building a Heap, and Using the Queue     sec:c07-build
%  Exports: thm:c07-floyd  fig:c07-buildfig  lst:c07-build  lst:c07-pq
%           ex:c07-build-run  prob:c07-topk  lst:c07-topk-soln
%           (practices: LC 215, LC 973, LC 451)
% =====================================================================
\section{Building a Heap, and Using the Queue}\label{sec:c07-build}

\begin{previouslybox}
\begin{enumerate}[nosep, leftmargin=1.6em]
  \item In \coderef{c07-heapops}, \code{push} sifts up a new leaf and \code{pop} sifts down from the root.
        Why does \emph{only that one path} need repair? (\defref{c07-sift} -- every untouched subtree is
        already a heap)
  \item A heap's height is $\lfloor\lg n\rfloor$ (\lemref{c07-height}). If you insert $n$ elements one at a
        time, what is the total cost -- and could a bulk build do better? ($\bigO{n\log n}$; yes, $\bigO{n}$,
        the point of this section)
  \item How many nodes of a complete tree on $n$ nodes sit at height exactly $h$ (roughly)? (about
        $n/2^{h+1}$ -- exponentially fewer as you go up)
  \item What does the standard \code{priority\_queue} default to -- min or max? (max-heap; \chref{c07})
\end{enumerate}
\end{previouslybox}

\begin{topicintro}
\secref{c07-heap} showed the two sifts but quietly assumed you \emph{have} a heap. Two questions remain,
and they are the difference between knowing the data structure and being able to use it under contest
pressure.

First, \textbf{build}: given a pile of $n$ unsorted values, turn them into a heap. Inserting one at a time
costs $\bigO{n\log n}$. Floyd's bottom-up heapify does it in $\bigO{n}$ (\thmref{c07-floyd}) -- because it
sifts each node down \emph{once}, and almost all nodes sit near the leaves where a sift moves only a
couple of steps. The proof is a geometric series and is one of the most re-quoted $\bigO{n}$-versus-$%
\bigO{n\log n}$ arguments in the subject; we do it on the running array $R$ in \exref{c07-build-run}.

Second, \textbf{use}: in practice you almost never hand-roll a heap -- you reach for
\code{std::priority\_queue}, and the whole skill is remembering its three quirks (it is a \emph{max}-heap
by default; \code{greater} flips it to min; the comparator's \code{bool} means ``\emph{is this one
smaller}'', so $\text{less} \Rightarrow$ the larger goes on top). \coderef{c07-pq} is the reference card.
Then the workhorse pattern: keep a \emph{size-$k$ heap} to track the top-$k$ / $k$-th-largest frontier in
$\bigO{\log k}$ per item without sorting -- cashed in on LeetCode 347, top $k$ frequent elements
(\probref{c07-topk}).
\textbf{Prerequisites.} \secref{c07-heap} (the sifts and the height bound) and the array layout of
\figref{c07-layout}.

\textbf{After this section you can}
\begin{itemize}[nosep]
  \item heapify an arbitrary array in $\bigO{n}$ bottom-up and prove the bound via the sum-of-heights
        series; \hfill\lv{4}
  \item drive \code{std::priority\_queue} -- min vs max, custom comparators, the \code{pair}/negation
        tricks -- without hand-rolling; \hfill\lv{3}
  \item maintain a size-$k$ heap for top-$k$ / $k$-th-largest and justify $\bigO{n\log k}$ over sorting. \hfill\lv{3}--\lv{4}
\end{itemize}
\end{topicintro}

\begin{recognitionbox}[Pattern card -- build once, then guard a frontier of size $k$]
\reckey{you are handed $n$ items to heapify \emph{at once} (so the $\bigO{n}$ build matters), or you must
track only the \emph{top $k$} / $k$-th extreme of a large or streaming input where $k \ll n$ and a full
sort is wasted work.}
\recuse{
  \item ``return the $k$ most frequent / $k$ closest / $k$-th largest'' -- a size-$k$ heap, \probref{c07-topk}
        and the practices;
  \item ``the library queue is your only container'' problems -- the \code{priority\_queue} idioms of
        \coderef{c07-pq};
  \item ``heapify this whole array, then answer many extract-min queries'' -- Floyd build,
        \coderef{c07-build}.
}
\recnotuse{
  \item the input is a \emph{stream} and you only ever need the single extreme -- a plain running
        max/min variable, no heap;
  \item you need the top $k$ \emph{and} want them in sorted order at the end with $k=\Theta(n)$ -- just
        sort, $\bigO{n\log n}$ beats $n$ heap-pops in constants;
  \item you need \emph{any} rank or the whole order maintained dynamically -- an ordered tree
        (\chref{c05}, \chref{c06}).
}
One-line test: \emph{sorting all $n$ to learn only $k$ things is the smell; a size-$k$ heap (or the
$\bigO{n}$ build) is the cure.}
\end{recognitionbox}

\subsection{Intuition}

Inserting $n$ elements one at a time pays $\bigO{\log n}$ for each because every insert may bubble all the
way up. But watch \emph{where} the work actually is in a complete tree: about half the nodes are leaves
(height $0$, they need no sift at all), a quarter are at height $1$ (one comparison), an eighth at height
$2$, and so on -- the count of nodes \emph{halves} every level you climb, while the cost of sifting from a
level \emph{doubles}. Those two exponentials cancel, leaving a convergent series that sums to linear. Floyd
exploits this by working \emph{bottom-up}: treat every node as the root of a tiny subtree whose children are
already heaps, and sift that one node down. By the time you reach the true root, its two subtrees are heaps
of size $\approx n/2$, so the single expensive sift-down at the top is the only $\bigO{\log n}$ move -- and
there is only one of them.

Run it on the array $R$ (indices $1..7$, values $5,3,8,1,4,2,7$). The last internal node is at index
$\lfloor 7/2\rfloor = 3$. We sift node $3$, then $2$, then $1$ -- and, as \exref{c07-build-run} walks, the
array settles into the very heap \figref{c07-heapR} shows, $(1,3,2,5,4,8,7)$, in a handful of swaps. No
insertion order, no recursion, one pass.

\begin{figure}[H]
\centering
\begin{tikzpicture}[diagram, scale=0.78, transform shape, x=13mm, y=9mm, level distance=13mm, sibling distance=20mm]
  % before
  \node[tnode] {$5$}
    child { node[tnode] {$3$} child { node[tleaf] {$1$} } child { node[tleaf] {$4$} } }
    child { node[tnode] {$8$} child { node[tleaf] {$2$} } child { node[tleaf] {$7$} } };
  \node[annot, font=\scriptsize\bfseries, text=warningred] at (0,1.2) {before: raw array $R$};
  \node[annot, align=left, font=\scriptsize] at (-4.6,-2.4)
    {sift node $3$ ($8$):\\ swap with $2$};
  \node[annot, align=left, font=\scriptsize] at (4.6,-2.4)
    {sift node $2$ ($3$):\\ swap with $1$};
  % after (to the right, small)
  \begin{scope}[xshift=8.2cm, x=9mm, y=7mm, level distance=9mm, sibling distance=15mm]
    \node[troot] {$1$}
      child { node[tnode] {$3$} child { node[tleaf] {$5$} } child { node[tleaf] {$4$} } }
      child { node[tnode] {$2$} child { node[tleaf] {$8$} } child { node[tleaf] {$7$} } };
    \node[annot, font=\scriptsize\bfseries, text=successgreen] at (0,0.95) {after: a heap};
  \end{scope}
  \draw[flow, ->] (2.6,-1.0) -- (5.4,-1.0) node[midway, above, font=\scriptsize] {Floyd};
\end{tikzpicture}
\caption{Floyd's bottom-up heapify turns the raw array $R=(5,3,8,1,4,2,7)$ into the min-heap
$(1,3,2,5,4,8,7)$ by sifting \emph{only the internal nodes}, from index $\lfloor n/2\rfloor$ up to the
root. Leaves are already (trivial) heaps; each sift-down merges two finished heaps plus one new value. The
expensive move happens once, at the root, which is why the total is $\bigO{n}$ (\thmref{c07-floyd},
\exref{c07-build-run}).}\label{fig:c07-buildfig}
\end{figure}

\subsection{The $\bigO{n}$ Build}

\begin{theorem}{Floyd's heapify is $\bigO{n}$}{c07-floyd}
Building a binary heap from $n$ arbitrary elements by the bottom-up procedure of \coderef{c07-build}
(sift-down each internal node once, from $\lfloor n/2\rfloor$ up to $1$) takes $\bigO{n}$ comparisons and
swaps, matching the $\bigO{n}$ information-theoretic lower bound up to constants and strictly beating the
$\bigO{n\log n}$ of repeated insertion.
\end{theorem}
\begin{proofbox}[Sum of heights converges]
A sift-down from a node costs $\bigO{\height(\text{node})}$, the height of the subtree rooted there. In a
complete tree on $n$ nodes, at most $\lceil n/2^{h+1}\rceil$ nodes have height $h$ (height $0$: the $\approx
n/2$ leaves; height $1$: $\approx n/4$; etc.). Hence the total work is at most
\[ \sum_{h\ge 0} \left\lceil \frac{n}{2^{h+1}} \right\rceil \cdot O(h)
   = O\!\left( n \sum_{h\ge 0} \frac{h}{2^{h+1}} \right). \]
The series $\sum_{h\ge0} h/2^{h+1} = \tfrac12\sum_{h\ge0} h/2^{h} = \tfrac12 \cdot 2 = 1$ converges (it is
the standard $\sum h x^{h} = x/(1-x)^2$ at $x=\tfrac12$). So the build is $O(n\cdot 1)=\bigO{n}$. The
naive $n$ separate inserts, by contrast, charge $\lg n$ to (almost) every node, giving $\bigO{n\log n}$ --
the gap is precisely that insert pays for the \emph{depth from the root}, while Floyd pays for the
\emph{height from the leaves}, and there are exponentially fewer short-sighted nodes to move. \end{proofbox}

\begin{lstlisting}[caption={Floyd's linear-time bottom-up heapify of \texttt{h[1..n]} (min-heap)}, label={lst:c07-build}]
void build(vector<int>& h, int n) {        // h[0] is an unused dummy; data in h[1..n]
    for (int i = n/2; i >= 1; --i) {       // every internal node, deepest-first
        int cur = i;
        while (true) {                    // one sift-DOWN; children are already heaps
            int s = cur, l = 2*cur, r = 2*cur + 1;
            if (l <= n && h[l] < h[s]) s = l;
            if (r <= n && h[r] < h[s]) s = r;
            if (s == cur) break;
            swap(h[cur], h[s]);
            cur = s;
        }
    }
}                          // total cost O(n) by thm:c07-floyd, not O(n log n)
\end{lstlisting}

\begin{complexitybox}
Time $\bigO{n}$ (\thmref{c07-floyd}); memory $\bigO{1}$ extra (in-place). Compare \coderef{c07-heapops}:
pushing the same $n$ values is $\bigO{n\log n}$. Both \emph{end} at a valid heap -- so the size-$k$ and
greedy patterns that follow are identical -- but if your heap starts life as a full array (the usual case:
you read $n$ numbers then want a queue), \code{build} it, do not insert it. Verdict: the single reason your
``heapify a vector'' step is not the bottleneck.
\end{complexitybox}

\subsection{The Library Priority Queue --- a Reference Card}

You rarely type \coderef{c07-build} in a contest; you write \code{priority\_queue}. Its four gotchas are
worth memorising once.

\begin{lstlisting}[caption={\code{std::priority\_queue} idioms: default max, \code{greater} for min, custom comparator, the size-k frontier}, label={lst:c07-pq}]
priority_queue<int> mx;                              // (1) MAX-heap by default
mx.push(5); int t = mx.top(); mx.pop();              // top() then pop(); no "pop returns"

priority_queue<int, vector<int>, greater<int>> mn;   // (2) MIN-heap via greater<>
// quick min-heap without writing a comparator: negate -- push(-x), read -(pq.top())

struct Cmp {   // (3) custom: operator() means "is a WORSE than b" => a below b
    bool operator()(const pair<int,int>& a, const pair<int,int>& b) const {
        if (a.second != b.second) return a.second < b.second; // max on .second
        return a.first > b.first;                              // tie-break min on .first
    }
};
priority_queue<pair<int,int>, vector<pair<int,int>>, Cmp> pq;

// (4) size-k frontier: a MIN-heap you never let exceed k
//     push each candidate, pop the root when size>k -> survivors = the k best
priority_queue<int, vector<int>, greater<int>> topk; // for k LARGEST: cap at k, min-root = kth
\end{lstlisting}

\begin{notebox}[Reading the comparator]
\code{priority\_queue} puts on \emph{top} the element that the comparator calls \textbf{largest}, and the
comparator \code{Cmp(a,b)} is ``$a$ should sit \emph{below} $b$''. So \code{less} ($a<b$) leaves the true
maximum on top (the default), and \code{greater} flips it to a min-heap. For a \code{pair} the default
orders lexicographically (by \code{.first} then \code{.second}); to prioritise by \code{.second} use a
comparator or store \code{\{second, first\}}. Confusing $\text{less}\Rightarrow\max$ with $\text{less}%
\Rightarrow\min$ is the number-one \code{priority\_queue} bug.
\end{notebox}

\subsection{Worked Example}

\begin{example}{Floyd build on $R=(5,3,8,1,4,2,7)$}{c07-build-run}
Array (indices $1..7$): $h=[5,3,8,1,4,2,7]$. Last internal index $\lfloor 7/2\rfloor = 3$; sift nodes
$3,2,1$.\\
\emph{$i=3$} ($h[3]{=}8$): children $h[6]{=}2$, $h[7]{=}7$; smallest is $2$; $8>2$ swap $\to$
$[5,3,2,1,4,8,7]$.\\
\emph{$i=2$} ($h[2]{=}3$): children $h[4]{=}1$, $h[5]{=}4$; smallest $1$; $3>1$ swap $\to$
$[5,1,2,3,4,8,7]$; node now $3$ at index $4$ has children $8,9$ out of range, stop.\\
\emph{$i=1$} ($h[1]{=}5$): children $h[2]{=}1$, $h[3]{=}2$; smallest $1$; $5>1$ swap $\to$
$[1,5,2,3,4,8,7]$; continue at index $2$: children $h[4]{=}3$, $h[5]{=}4$, smallest $3$; $5>3$ swap $\to$
$[1,3,2,5,4,8,7]$; index $4$ children out of range, stop.\\
Final heap $[1,3,2,5,4,8,7]$ -- identical to \figref{c07-heapR} and to the insert-order result of
\exref{c07-sift-run}, reached in five swaps. Verify every parent $\le$ children: $1\le3,2$;\ $3\le5,4$;\
$2\le8,7$. \checkmark
\end{example}
\begin{tracebox}[Swap ledger for \exref{c07-build-run} -- only indices 3,2,1 sifted]
i=3: swap(3,6) -> 5 3 2 1 4 8 7
i=2: swap(2,4) -> 5 1 2 3 4 8 7
i=1: swap(1,2) -> 1 5 2 3 4 8 7 ; then swap(2,4) -> 1 3 2 5 4 8 7
total swaps = 4 (plus the one at i=3), cost = O(height) each: i=3 height 1, i=2 height 1, i=1 height 2.
\end{tracebox}

\subsection{Variations and Twists}

\begin{itemize}[nosep]
  \item \textbf{Floyd vs.\ inserts}: both yield \emph{a} valid heap, generally \emph{different} arrays (any
        of several arrays is a legal heap on the same multiset). Never claim ``the'' heap of an array --
        only ``a'' heap -- unless you also fix the build procedure. \hfill\lv{3}
  \item \textbf{D-ary heaps}: trade height for branching. A $d$-ary heap has height $\log_d n$ (insert
        cheaper) but sift-down compares $d$ children (pop dearer): $\bigO{d\log_d n}$ per pop. Large $d$
        wins on cache locality; the priority-queue-with-decreases-key of Dijkstra likes small $d$. Ideas
        only here -- see also \secref{c07-trees}. \hfill\lv{4}
  \item \textbf{Heapsort}: pop the root $n$ times from a max-heap, writing each winner into the vacated
        tail -- an in-place $\bigO{n\log n}$, $\bigO{1}$-extra sort. Its cousin ``$k$ passes of selection''
        is exactly the size-$k$ frontier when you only want $k$ outputs. \hfill\lv{4}
  \item \textbf{build + lazy delete}: Floyd the array, then to ``remove'' an arbitrary element mark it dead
        and pop stale marked tops on demand -- the trick the Summit (\probref{c07-summit}) uses to get
        $\bigO{\log n}$ ``delete'' of a non-root without a position map. \hfill\lv{4}
\end{itemize}

\subsection{Common Pitfalls}

\begin{warningbox}[Trap 1 -- assuming \code{priority\_queue} is a min-heap]
The default \code{priority\_queue<T>} is a \textbf{max}-heap: \code{top()} gives the \emph{largest}. A
``find the $k$-th \emph{largest} with a size-$k$ heap'' needs a \emph{min}-heap (so the small ones get
evicted and the root is the $k$-th largest). Forgetting the \code{greater<int>} (or the negation) silently
keeps the $k$ \emph{smallest} instead -- a wrong answer that still runs in $\bigO{n\log k}$, so nothing
crashes to warn you.
\end{warningbox}
\begin{lstlisting}[caption={\textcolor{warningred}{WRONG}: default (max) heap for ``top k largest'' keeps the k smallest}, language=C++, basicstyle=\ttfamily\small]
priority_queue<int> pq;                       // max-heap!
for (int x : a){ pq.push(x); if (pq.size()>k) pq.pop(); }  // pops the LARGEST each time
// survivors are the k SMALLEST -- the opposite of the intent
\end{lstlisting}
\begin{lstlisting}[caption={\textcolor{successgreen}{RIGHT}: a size-k MIN-heap keeps the k largest}, language=C++, basicstyle=\ttfamily\small]
priority_queue<int, vector<int>, greater<int>> pq;   // min-heap
for (int x : a){ pq.push(x); if (pq.size()>k) pq.pop(); }  // evict the smallest
// survivors are the k largest; pq.top() is the k-th largest
\end{lstlisting}

\begin{warningbox}[Trap 2 -- building by insertion when you already have the array]
Calling \code{push} $n$ times is $\bigO{n\log n}$; if you are handed the whole array up front, that is a
needless extra $\log n$ factor on a problem whose real bottleneck was meant to be elsewhere. Use the
$\bigO{n}$ Floyd build (\thmref{c07-floyd}), or in C++ the constructor
\code{priority\_queue<int>(make\_move\_iterator...)} route / \code{make\_heap} on a vector. For $n=10^{6}$
the difference is measurable.
\end{warningbox}

\begin{warningbox}[Trap 3 -- comparing by \code{.first} when you meant \code{.second}]
Storing \code{\{count, value\}} and wanting the \emph{max count} works with the default max-heap (it
orders by \code{.first} first). Storing \code{\{value, count\}} and expecting it to sort by count does
\emph{not} -- the lexicographic order reads \code{value} first. Put the \emph{priority} in the first field
of the \code{pair}, or supply a comparator. This is the second-most-common \code{priority\_queue} error
after the min/max flip.
\end{warningbox}

\subsection{Problem Lab}

LeetCode 347 asks for the $k$ most frequent elements -- the archetypal ``$k \ll n$, do not sort everything''
problem. Two passes and a \emph{size-$k$ min-heap} over frequency pairs do it in $\bigO{n\log k}$; we also
note the $\bigO{n}$ bucket-sort alternative that sidesteps the heap when values are bounded.

% ---------------------------------------------------------------------------
\solvedproblem{c07-topk}{Top K Frequent Elements}{LeetCode 347}
  {https://leetcode.com/problems/top-k-frequent-elements/}{LC Medium · L3}
  {hash-map frequencies, size-k min-heap frontier, bucket-sort alternative}

\begin{problemstatement}[Top K Frequent Elements]
Given an integer array \texttt{nums} and an integer $k$, return the $k$ \textbf{most frequent} elements. You
may return the answer in \textbf{any order}. Constraints: $1 \le \texttt{nums.length} \le 10^{5}$, $k$ is in
the range of the number of distinct values, and it is guaranteed the answer is unique.
\end{problemstatement}

\begin{attemptgate}[Attempt it yourself first -- \lv{3}, give yourself 15 minutes]
Counting frequencies is trivial; the question is how to get the top $k$ \emph{without} fully sorting the
distinct values. Keep only $k$ candidates in a heap whose root is the \emph{least} frequent of the current
$k$.
\peek{After building \code{freq[value]}, what does a min-heap keyed by frequency, capped at size $k$,
guarantee about what survives?}
\nudge{Push every \code{\{freq, value\}} pair onto a \emph{min}-heap; when the heap exceeds $k$ elements,
pop its root (the least frequent of those seen). The $k$ survivors are the $k$ most frequent.}
\fullhint{Because the heap is a min-heap on frequency, its \emph{root} is always the current \emph{weakest}
of the top-$k$ set, so evicting the root is safe -- a better candidate can only have higher frequency and
will displace the weakest. The $k$ values left in the heap are the answer. For \texttt{nums=[1,1,1,2,2,3]},
$k=2$: freqs are $1{:}3, 2{:}2, 3{:}1$; survivors $1$ and $2$. Answer $[1,2]$.}
\end{attemptgate}

\solstep{Understanding the Problem}
We need the $k$ values with the highest occurrence counts, order of output irrelevant. A full sort of all
distinct values by frequency is $\bigO{m\log m}$ where $m$ is the number of distinct values; we only need
the top $k$, so a size-$k$ heap is $\bigO{m\log k}$, and when $k \ll m$ that is a real saving. The unique-answer
guarantee means no frequency ties straddle the $k$-boundary, so ``which $k$ survive'' is well-defined.

\solstep{Key Observations}
\begin{observation}
Frequencies come from one hash-map pass, $\bigO{n}$; the heap only sees the $m$ distinct \code{\{freq,
value\}} pairs.
\end{observation}
\begin{observation}
A \emph{min}-heap on frequency capped at size $k$ evicts the \emph{weakest} current candidate, so it
always holds the $k$ strongest seen so far -- the root is the $k$-th-strongest, our ``frontier''
(\coderef{c07-pq}, idiom 4).
\end{observation}
\begin{observation}
Since $1 \le \text{freq} \le n$, an $\bigO{n}$ \emph{bucket} array (index $=$ frequency) reads off the top
$k$ in $\bigO{n}$ worst-case with no heap at all -- the preferred solution when $n$ is small enough to
bucket. We give the heap version because the pattern generalises to unbounded/weighted top-$k$.
\end{observation}

\solstep{Think It Through With Me}
\begin{thinkbox}
\deadend{Sort all distinct values by frequency, take the first $k$.}
{$\bigO{m\log m}$; correct but does $\log m$ work per value to learn only $k$ things. Fine here, wasteful
as a habit, and it hides the size-$k$ frontier idea you need elsewhere.}
\deadend{Use a \emph{max}-heap and pop $k$ times.}
{Works ($\bigO{n + k\log m}$ after an $\bigO{n}$ build) but pays a full $\bigO{m}$ heap; if $m\gg k$ the
size-$k$ min-heap is cheaper in memory and the constant is smaller. Choose deliberately, not by reflex.}
\aha{Cap a \emph{min}-heap at $k$: every insert that pushes size to $k{+}1$ evicts the root, which is
exactly the current least-frequent survivor. $\bigO{m\log k}$, and the root is the $k$-th most frequent.}
\end{thinkbox}

\solstep{Solution}
Build \code{freq} with a hash map. Iterate the distinct values, pushing \code{\{frequency, value\}} onto a
\textbf{min}-heap keyed by frequency; whenever the heap exceeds size $k$, pop the root. At the end the heap
holds the $k$ most frequent values; drain their \code{.second} fields into the answer. Uniqueness of the
answer removes tie ambiguity at the boundary.

\solstep{Proof of Correctness}
\begin{proofbox}[The capped min-heap holds the $k$ best seen so far]
Invariant: after processing any prefix of the distinct values, the heap contains the $\min(k, t)$ highest-
frequency values among the $t$ seen, and if $t>k$ its root is the least frequent of those. Base $t\le k$:
nothing is evicted, so all $t$ are held -- true. Step: on the $(t{+}1)$-th pair we push it then, if size $k{+}1$,
pop the root (smallest frequency in the heap). The popped element is $\le$ every remaining heap element, so
removing it cannot remove a member of the top-$k$ of the now $t{+}1$ seen values: the top $k$ of $t{+}1$
values are exactly the previous top $k$ with the new pair possibly inserted and the weakest of the union
dropped, which is what push-then-pop does. Hence the invariant holds. At termination the heap is the
top-$k$ of all $m$ distinct values -- the answer. \end{proofbox}

\solstep{Implementation}
\begin{lstlisting}[caption={Top K Frequent Elements -- size-k min-heap over frequency pairs (LeetCode 347)}, label={lst:c07-topk-soln}]
class Solution {
public:
    vector<int> topKFrequent(vector<int>& nums, int k) {
        unordered_map<int,int> freq;
        for (int x : nums) ++freq[x];                 // O(n) counting pass
        priority_queue<pair<int,int>,                  // min-heap keyed by .first=freq
                       vector<pair<int,int>>, greater<>> pq;
        for (auto& p : freq) {
            pq.push({p.second, p.first});              // {frequency, value}
            if ((int)pq.size() > k) pq.pop();          // evict least frequent
        }
        vector<int> ans;
        while (!pq.empty()) { ans.push_back(pq.top().second); pq.pop(); }
        return ans;                                    // any order
    }
};
\end{lstlisting}

\solstep{Trace}
\begin{tracebox}[Official sample: \texttt{nums=[1,1,1,2,2,3]}, \texttt{k=2}]
freq $= \{1{:}3,\, 2{:}2,\, 3{:}1\}$. Push pairs \code{\{freq,value\}}: $(3,1)$ size1; $(2,2)$ size2;
$(1,3)$ size3$>2$ so pop the root $=$ least frequent $=(1,3)$. Survivors $\{(2,2),(3,1)\}$, i.e.\ values
$2$ and $1$. Drain $=$ an order of $[1,2]$. \checkmark\ matches the official output $[1,2]$ (any order).
\end{tracebox}

\solstep{Complexity}
\begin{complexitybox}
Time $\bigO{n + m\log k}$ where $m$ is the number of distinct values ($m \le n$): one map pass plus one
push and occasional pop per distinct value, each heap op $\bigO{\log k}$. Memory $\bigO{m}$ for the map and
the size-$k$ heap. The $\bigO{n}$ \emph{bucket} variant (an array of lists indexed by frequency, scanned
downward) drops the $\log k$ but costs $\bigO{n}$ memory and only applies when frequencies are integers
$\le n$ -- know both; the heap one is the transferable top-$k$ pattern (\probref{c07-summit} reuses the
frontier idea with a different key). Verdict: \ac.
\end{complexitybox}

\solstep{Pitfalls and Variants}
\begin{warningbox}
Ways to fall: (i) keying the \emph{max}-heap by frequency and popping $k$ times but forgetting you must
\emph{build} it in $\bigO{n}$ (Floyd) or it is secretly $\bigO{m\log m}$; (ii) using \code{greater<>} and
then reading the \emph{wrong} field -- the root is the least frequent, so the answer is the \emph{values}
still in the heap, not their frequencies; (iii) assuming output order matters -- it does not here, but for
\emph{Sort Characters By Frequency} (LC 451, a practice) it does, so there drain a max-heap.
\emph{Variant:} LC 692 (top-$k$ words) adds a lexicographic tie-break -- put both keys in the \code{pair}
and let the comparator order by \code{\{freq desc, word asc\}}.
\end{warningbox}

\begin{insightbox}[Key takeaways]
\begin{enumerate}[nosep, leftmargin=1.6em]
  \item \textbf{Build in $\bigO{n}$} by sifting internal nodes bottom-up (\thmref{c07-floyd},
        \coderef{c07-build}); inserting is the $\bigO{n\log n}$ default you should beat when you hold the
        whole array.
  \item \textbf{\code{priority\_queue} is a max-heap}; \code{greater} makes it a min-heap; a \code{pair}'s
        \emph{first} field is the priority. Four gotchas, one card (\coderef{c07-pq}).
  \item \textbf{Top-$k$ is a size-$k$ heap}, not a sort: cap the frontier, evict the root, answer in
        $\bigO{n\log k}$ -- \probref{c07-topk} and the whole practice set.
\end{enumerate}
\end{insightbox}

\begin{hardeningbox}
\begin{itemize}[nosep, leftmargin=1.4em]
  \item \textbf{Rebuild $R$} with \coderef{c07-build} by hand and confirm you land on $(1,3,2,5,4,8,7)$
        exactly as \exref{c07-build-run} does.
  \item \textbf{Bucket version} of \probref{c07-topk}: allocate \code{vector<vector<int>> bucket(n+1)},
        place each value at index $=$ its frequency, then scan downward collecting $k$; argue it is $\bigO{n}$.
  \item \textbf{Comparator drill}: rewrite \coderef{c07-topk-soln}'s heap as a \emph{max}-heap and pop $k$
        times -- same answer, different cost -- and state when each wins.
\end{itemize}
\end{hardeningbox}

\subsection{Practice Problems}

\begin{stuckbox}
\begin{itemize}[nosep, leftmargin=1.4em]
  \item Did you choose min- or max-heap to match \emph{which} end you evict in the size-$k$ frontier?
  \item Did you put the priority in the \code{pair}'s \emph{first} field (or write a comparator)?
  \item When you were handed a full array, did you \emph{build} the heap in $\bigO{n}$ rather than insert?
\end{itemize}
\end{stuckbox}

\begin{practiceset}
  \practice{Kth Largest Element in an Array}{LeetCode 215}
    {https://leetcode.com/problems/kth-largest-element-in-an-array/}{LC Medium · L3}
    {Return the $k$-th largest element (sorted order of the array, not distinct). Two heap-grade answers:
    a size-$k$ \emph{min}-heap leaves the $k$ largest with the $k$-th largest at the root ($\bigO{n\log k}$);
    or build one \emph{max}-heap in $\bigO{n}$ (Floyd, \thmref{c07-floyd}) and pop $k$ times
    ($\bigO{n + k\log n}$). Note the array may contain duplicates -- the $k$-th \emph{largest by position},
    so count carefully. \textbf{Also} know the expected-$\bigO{n}$ quickselect (partition on a random pivot
    and recurse on one side) -- the interview favourite when $\bigO{\log} $ is unwanted.}
  \practice{K Closest Points to Origin}{LeetCode 973}
    {https://leetcode.com/problems/k-closest-points-to-origin/}{LC Medium · L3}
    {Return the $k$ points closest to $(0,0)$ (distance $\sqrt{x^2+y^2}$; compare squared distance to avoid
    the sqrt and any overflow -- use \code{long long}). Exactly \probref{c07-topk} with a computed key: keep
    a size-$k$ \emph{max}-heap keyed by squared distance so you evict the \emph{farthest} current candidate,
    leaving the $k$ nearest ($\bigO{n\log k}$). Order of output does not matter.}
  \practice{Sort Characters By Frequency}{LeetCode 451}
    {https://leetcode.com/problems/sort-characters-by-frequency/}{LC Medium · L3}
    {Reorder a string so characters appear in descending frequency (ties any order). Count frequencies, then
    drain a \emph{max}-heap keyed by count -- unlike \probref{c07-topk} you want \emph{all} characters in
    order, so pop the whole heap, appending each char its count times. $\bigO{n + c\log c}$ for $c$ distinct
    chars; a bucket array gives $\bigO{n}$. This is the ``heap as a priority-driven output iterator''
    pattern.}
\end{practiceset}

\begin{recallbox}[10-second refresher -- 7.2 Build + Use]
Floyd heapifies an array in $\bigO{n}$ by sifting internal nodes from $\lfloor n/2\rfloor$ down-to-up
(\thmref{c07-floyd}, \coderef{c07-build}); the sum-of-heights series $\sum h/2^{h+1}$ converges, which is why
it beats $\bigO{n\log n}$ insertion. \code{priority\_queue} is a \emph{max}-heap by default, use
\code{greater} for min, a \code{pair}'s first field is its priority (\coderef{c07-pq}). Top-$k$ /
$k$-th-largest is a \emph{size-$k$} heap ($\bigO{n\log k}$), not a sort -- \probref{c07-topk}.
\selftest{
  \item Why does Floyd cost $\bigO{n}$ while $n$ inserts cost $\bigO{n\log n}$? State the series.
  \item What is the default flavour of \code{priority\_queue} and how do you get a min-heap two ways?
  \item For top-$k$ \emph{largest}, is the capped heap a min- or max-heap, and what is its root?
  \item In a \code{\{a,b\}} heap element, which field decides priority, and how do you prioritise the other?
  \item When does bucket sort beat the heap for top-$k$, and at what memory cost?
}
\end{recallbox}
